\makesectionpage{Unüberwachtes Lernen}

\makesubsectionpage{Abstandserhaltende Projektionen}

\begin{frame}{Abstandserhaltende Projektionen}
\begin{columns}
\column{0.55\textwidth}
\begin{itemize}
  \item Oft möchte man \textbf{hochdimensionale} Daten visuell analysieren.
  \item Interessant ist vor allem die \textbf{Verteilung} der Datenpunkte im
        Raum.
  \item Eine \textbf{abstandserhaltende Projektion} bildet Punkte aus einem
        Ausgangsraum in einen Zielraum \textbf{geringerer Dimension} ab.
  \item Dabei bleiben die \textbf{Abstände} in etwa erhalten: nahe Punkte
        bleiben nah, ferne bleiben fern.
\end{itemize}

\column{0.35\textwidth}
\centering
\resizebox{\linewidth}{!}{%
\begin{tikzpicture}[>=Stealth,
  pa/.style={circle, fill=SecondaryColor, draw=SecondaryTitle,
             very thick, minimum size=5pt, inner sep=0pt},
  pb/.style={circle, fill=AccentColor, draw=AccentTitle,
             very thick, minimum size=5pt, inner sep=0pt}
]
% Ausgangsraum (3D angedeutet)
\draw[gray!50] (0,2.6) -- (3,2.6) -- (3.8,3.3) -- (0.8,3.3) -- cycle;
\draw[gray!50] (0,2.6) -- (0,4.0) -- (0.8,4.7) -- (0.8,3.3);
\draw[gray!50] (3,2.6) -- (3,4.0) -- (3.8,4.7) -- (3.8,3.3);
% Cluster A (blau, eng)
\node[pa] at (0.8,3.6) {};
\node[pa] at (1.2,3.8) {};
\node[pa] at (1.2,3.3) {};
% Cluster B (orange, entfernt)
\node[pb] at (2.8,3.9) {};
\node[pb] at (2.6,3.6) {};
\node[pb] at (2.9,4.1) {};
\node[pb] at (2.4,4.0) {};
\node[font=\scriptsize\bfseries] at (1.9,4.5) {Ausgangsraum};

% Projektionspfeil
\draw[->, thick] (1.9,2.4) -- (1.9,1.7);
\node[font=\scriptsize, right] at (2.0,2.05) {Projektion};

% Zielraum (2D)
\node[font=\scriptsize\bfseries] at (1.9,1.4) {Zielraum};
\draw[thick] (0.2,-0.2) rectangle (3.4,1.2);
% Cluster A (blau, eng)
\node[pa] at (0.8,0.4) {};
\node[pa] at (1.1,0.6) {};
\node[pa] at (0.9,0.2) {};
% Cluster B (orange, entfernt)
\node[pb] at (2.7,0.7) {};
\node[pb] at (3.0,0.4) {};
\node[pb] at (2.8,0.9) {};
\node[pb] at (2.5,0.5) {};
\end{tikzpicture}}
\end{columns}
\end{frame}

\begin{frame}{Abstandserhaltende Projektionen: Beispiel}
\begin{columns}
\column{0.47\textwidth}
\begin{itemize}
  \item Daten sollen aus der \textbf{Ebene} auf eine \textbf{Gerade}
        projiziert werden.
  \item Das ist eine \textbf{Dimensionsreduktion} von \textbf{2 auf 1}.
  \item Ziel ist es, die \textbf{Nachbarschaftsbeziehungen} der Punkte zu
        erhalten.
  \item Getrennte Cluster bleiben auch nach der Projektion \textbf{getrennt}.
\end{itemize}

\column{0.48\textwidth}
\centering
\resizebox{\linewidth}{!}{%
\begin{tikzpicture}[>=Stealth,
  gr/.style={circle, fill=AccentColor, draw=AccentTitle,
             very thick, minimum size=6pt, inner sep=0pt},
  ma/.style={circle, fill=AlertColor, draw=AlertTitle,
             very thick, minimum size=6pt, inner sep=0pt},
  bl/.style={circle, fill=SecondaryColor, draw=SecondaryTitle,
             very thick, minimum size=6pt, inner sep=0pt}
]
% Achsen (Ebene)
\draw[thick, SecondaryTitle] (0,5.5) -- (0,2.0) -- (6.5,2.0);

% gruenes Cluster (oben links)
\node[gr] at (1.55,4.75) {};
\node[gr] at (1.25,4.35) {};
\node[gr] at (1.8,4.45) {};
\node[gr] at (1.5,4.0) {};
% magenta Cluster (mitte, unterhalb gruen)
\node[ma] at (1.9,3.1) {};
\node[ma] at (1.6,2.7) {};
\node[ma] at (2.2,2.6) {};
\node[ma] at (2.0,2.3) {};
% blaues Cluster (oben rechts)
\node[bl] at (4.65,4.95) {};
\node[bl] at (4.3,4.5) {};
\node[bl] at (4.95,4.5) {};
\node[bl] at (4.6,4.1) {};

% Zielgerade (1D)
\draw[thick] (0,0) -- (6.5,0);
% projizierte Cluster: alle drei getrennt
\foreach \x in {0.6,1.0,1.3,1.7 } \node[gr] at (\x,0.25) {};
\foreach \x in {2.5,2.8,3.2,3.6} \node[ma] at (\x,0.25) {};
\foreach \x in {4.8, 5.1,5.6, 6.0} \node[bl] at (\x,0.25) {};

% Projektionspfeile
\draw[->, line width=1pt] (1.3,3.9) to[bend right=20] (1.2,0.55);
\draw[->, line width=1pt] (2.5,2.4) to[bend left=15] (3.2,0.55);
\draw[->, line width=1pt] (5.0,3.9) -- (5.2,0.55);
\end{tikzpicture}}
\end{columns}
\source{\cite{seemann2024mlaz}}
\end{frame}


\begin{frame}{Abstandserhaltende Projektionen: Problem}
\begin{columns}
\column{0.47\textwidth}
\begin{itemize}
  \item Eine \textbf{einfache Projektion} auf eine der Koordinatenachsen ist
        nicht abstandserhaltend.
  \item Cluster, die in der Ebene \textbf{getrennt} sind, können dabei
        \textbf{überlappen}.
  \item Die \textbf{Nachbarschaftsinformation} geht so verloren.
  \item Hier fallen zwei Cluster
        zusammen.
\end{itemize}

\column{0.48\textwidth}
\centering
\resizebox{\linewidth}{!}{%
\begin{tikzpicture}[>=Stealth,
   gr/.style={circle, fill=AccentColor, draw=AccentTitle,
             very thick, minimum size=6pt, inner sep=0pt},
  ma/.style={circle, fill=AlertColor, draw=AlertTitle,
             very thick, minimum size=6pt, inner sep=0pt},
  bl/.style={circle, fill=SecondaryColor, draw=SecondaryTitle,
             very thick, minimum size=6pt, inner sep=0pt}
]
% Achsen (Ebene)
\draw[thick, SecondaryTitle] (0,5.5) -- (0,2.0) -- (6.5,2.0);

% gruenes Cluster (oben links)
% gruenes Cluster (oben links)
\node[gr] at (1.55,4.75) {};
\node[gr] at (1.25,4.35) {};
\node[gr] at (1.8,4.45) {};
\node[gr] at (1.5,4.0) {};
% magenta Cluster (mitte, unterhalb gruen)
\node[ma] at (1.9,3.1) {};
\node[ma] at (1.6,2.7) {};
\node[ma] at (2.2,2.6) {};
\node[ma] at (2.0,2.3) {};
% blaues Cluster (oben rechts)
\node[bl] at (4.65,4.95) {};
\node[bl] at (4.3,4.5) {};
\node[bl] at (4.95,4.5) {};
\node[bl] at (4.6,4.1) {};

% Zielgerade (1D)
\draw[thick] (0,0) -- (6.5,0);
% projizierte Cluster: gruen und magenta ueberlappen
\foreach \x in {1.25,1.55,1.8, 1.5} \node[gr] at (\x,0.25) {};
\foreach \x in {1.6,1.9,2.2, 2.0} \node[ma] at (\x,0.25) {};
\foreach \x in {4.3,4.65,4.95, 4.6} \node[bl] at (\x,0.25) {};

% Projektionspfeile (senkrecht)
\draw[->, line width=1pt] (1.8,2.3) -- (1.8,0.55);
\draw[->, line width=1pt] (4.6,3.8) -- (4.6,0.55);
\end{tikzpicture}}
\end{columns}
\source{\cite{seemann2024mlaz}}
\end{frame}

\begin{frame}{MDS: Grundidee}
\footnotesize
\begin{columns}
\column{0.55\textwidth}
\begin{itemize}
  \item \textbf{MDS} (Multidimensional Scaling) ist eine
        \textbf{abstandserhaltende} Projektion.
  \item \textbf{Eingabe:} paarweise \textbf{Unähnlichkeiten} $\delta_{ij}$
        zwischen $n$ Objekten.
  \item \textbf{Ziel:} eine Konfiguration $\mathbf{x}_1, \dots, \mathbf{x}_n$
        in \textbf{niedriger} Dimension (meist $\mathbb{R}^2$).
  \item Die Abstände $d_{ij}(X) = \lVert \mathbf{x}_i - \mathbf{x}_j \rVert$
        sollen die $\delta_{ij}$ möglichst gut \textbf{annähern}.
\end{itemize}

\begin{blockC1}{\footnotesize Kernidee}
Statt der \textbf{Koordinaten} der Ausgangspunkte werden nur ihre
\textbf{Abstände} zueinander benötigt.
\end{blockC1}

\column{0.40\textwidth}
\begin{blockC2}{\footnotesize Konfiguration}
Die Konfiguration $X$ wird als \textbf{Matrix} dargestellt:
\[
X=
\begin{pmatrix}
\mathbf{x}_1^T\\
\mathbf{x}_2^T\\
\vdots\\
\mathbf{x}_n^T
\end{pmatrix}
\in \mathbb{R}^{n\times p}.
\]

\medskip

\textbf{Jede Zeile} entspricht einem der projizierten Punkte $\mathbf{x}_1, \dots, \mathbf{x}_n \in \mathbb{R}^p$.

\medskip

Wollen wir $n$ Punkte in einen 2D-Zielraum projizieren, dass ist $X$ ein $n\times 2$-Matrix
\end{blockC2}


\end{columns}
\source{\cite{deLeeuw1977}}
\end{frame}

\begin{frame}{MDS: Die Stress-Funktion}
  \small
\begin{columns}
\column{0.55\textwidth}
Die Güte einer Konfiguration $X$ misst man über die \textbf{Stress}-Funktion $\sigma: \mathbb{R}^{n\times p} \to \mathbb{R}$, definiert durch 
\[
  \sigma(X) = \sum_{i<j} w_{ij}\bigl(\delta_{ij} - d_{ij}(X)\bigr)^2
\]
für alle $X\in \mathbb{R}^{n\times p}$. 

\begin{itemize}
  \item $\delta_{ij}$: gewünschte \textbf{Zielabstände}. Werden initial berechnet und sind dann fix.
  \item $d_{ij}(X)$: tatsächliche \textbf{Abstände} in der Projektion. Ändern sich, wenn die Zielpunkte während der Optimierung verschoben werden. 
  \item $w_{ij} \geq 0$: optionale \textbf{Gewichte}.
\end{itemize}

\column{0.40\textwidth}


\begin{blockDef}{\small Ziel}
  Gesucht ist das $X$, das den Stress \textbf{minimiert}:
\[ \min_X \; \sigma(X)\]

\medskip
Kleiner Stress bedeutet: Die \textbf{Abstände} in der Projektion passen gut
zu den \textbf{Zielabständen}.

\medskip
$\Rightarrow$ Wir haben wieder ein \textbf{Optimierungsproblem}.
\end{blockDef}
\end{columns}
\source{\cite{deLeeuw1977}}
\end{frame}

\begin{frame}{MDS: Warum ist das schwierig?}
  \footnotesize
\begin{columns}
\column{0.60\textwidth}
Ein Optimierungsproblem heißt \textbf{konvex}, wenn die Funktion nur
\textbf{ein einziges} und damit globales  Minimum besitzt. Solche Probleme lassen sich zuverlässig lösen.

\vspace{0.4em}
\begin{itemize}
  \item Die Stress-Funktion ist \textbf{nicht konvex}: Sie hat \textbf{viele
        lokale Minima}.
  \item Der Grund: $d_{ij}(X) = \lVert \mathbf{x}_i - \mathbf{x}_j \rVert
  = \sqrt{\sum_{a=1}^{p} (x_{ia} - x_{ja})^2}$
        enthält eine \textbf{Wurzel} und ist selbst nicht konvex.
  \item Der \textbf{Gradientenabstieg} läuft in das \textbf{nächstgelegene}
        lokale Minimum, je nach Startpunkt ein anderes.
  \item Das globale Minimum wird so \textbf{nicht garantiert} gefunden.
\end{itemize}

\column{0.35\textwidth}
\begin{blockC2}{\footnotesize Konsequenz}
Statt reinem Gradientenabstieg nutzt man \textbf{SMACOF}, das den Stress in
jedem Schritt garantiert \textbf{senkt}.
\end{blockC2}
\end{columns}
\source{\cite{deLeeuw1977}}
\note{Eine Funktion heißt \textbf{konvex}, wenn zwischen zwei beliebigen Punkten auf ihrem Graphen die Verbindungslinie \textbf{oberhalb} des Graphen verläuft.\\
Konvexe Funktionen haben die Form einer Schüssel. Sie haben ein eindeutiges Minimum.
}
\end{frame}

\begin{frame}{SMACOF: Majorisierung}
\begin{columns}
\column{0.60\textwidth}
\textbf{SMACOF} steht für \emph{Scaling by MAjorizing a COmplicated
Function}.

\vspace{0.4em}
\begin{itemize}
  \item Statt den Stress $\sigma(X)$ direkt zu minimieren, sucht man eine
        \textbf{einfachere} Hilfsfunktion, die \textbf{Majorante}.
  \item Dabei ist $Z = X^{(k)}$ die \textbf{aktuelle} Konfiguration des
        letzten Schritts.
  \item Die Majorante liegt \textbf{überall oberhalb} von $\sigma$ und
        \textbf{berührt} es im aktuellen Punkt $Z$.
  \item Sie ist \textbf{quadratisch} und daher leicht zu minimieren.
\end{itemize}

\column{0.35\textwidth}
\begin{blockC1}{Prinzip}
Minimiert man die Majorante und setzt das Ergebnis als neues $Z$ ein,
\textbf{sinkt} der Stress garantiert.
\end{blockC1}
\end{columns}
\source{\cite{deLeeuw1977}}
\end{frame}

%%%%%%%%%%%%%%%%%%%%%%%%%%%%%%%%%%%%%%%%%%%%%
\begin{frame}{SMACOF: Die Spur als Werkzeug}
\begin{columns}
\column{0.55\textwidth}
Die \textbf{Spur} $\operatorname{tr}(A)$ einer quadratischen Matrix ist die
Summe ihrer \textbf{Diagonalelemente}:
\[
  \operatorname{tr}(A) = \sum_{i} a_{ii}
\]

\begin{itemize}
  \item Sie erlaubt es, \textbf{Summen von Quadraten} kompakt als
        Matrixausdruck zu schreiben.
  \item Die Abstände lassen sich damit darstellen als
        $\sum_{i<j} w_{ij}\, d_{ij}^2(X) = \operatorname{tr}(X^T V X)$.
  \item So werden die vielen \textbf{paarweisen} Terme in \textbf{eine}
        Matrixformel gebündelt.
\end{itemize}

\column{0.40\textwidth}
\begin{blockC2}{Matrix $V$}
\[
  v_{ij} =
  \begin{cases}
    -w_{ij} & i \neq j\\[2pt]
    \sum_{k \neq i} w_{ik} & i = j
  \end{cases}
\]
\end{blockC2}

\begin{blockC1}{Nutzen}
Die Spur-Schreibweise macht den Stress \textbf{ableitbar} und offenbart
seine \textbf{quadratische} Struktur in $X$.
\end{blockC1}
\end{columns}
\source{\cite{deLeeuw1977}}
\end{frame}

\begin{frame}{SMACOF: Herleitung der Spur-Identität}
\begin{columns}
\column{0.60\textwidth}
Wir zeigen
$\sum_{i<j} w_{ij}\, d_{ij}^2(X) = \operatorname{tr}(X^T V X)$.

\vspace{0.3em}
Mit $d_{ij}^2(X) = \lVert \mathbf{x}_i - \mathbf{x}_j \rVert^2$ und der
Definition von $V$ rechnet man:
\[
  \begin{aligned}
    \operatorname{tr}(X^T V X)
      &= \sum_{i,j} v_{ij}\,\mathbf{x}_i \cdot \mathbf{x}_j\\
      &= \sum_{i} \Bigl(\sum_{k\neq i} w_{ik}\Bigr)\mathbf{x}_i \cdot \mathbf{x}_i
         - \sum_{i\neq j} w_{ij}\,\mathbf{x}_i \cdot \mathbf{x}_j\\
      &= \sum_{i<j} w_{ij}
         \bigl(\mathbf{x}_i \cdot \mathbf{x}_i - 2\,\mathbf{x}_i \cdot \mathbf{x}_j
         + \mathbf{x}_j \cdot \mathbf{x}_j\bigr)\\
      &= \sum_{i<j} w_{ij}\,\lVert \mathbf{x}_i - \mathbf{x}_j \rVert^2.
  \end{aligned}
\]

\column{0.35\textwidth}
\begin{block}{Kernschritt}
Das \textbf{Diagonalelement} $\sum_{k\neq i} w_{ik}$ liefert genau die
Terme $\mathbf{x}_i \cdot \mathbf{x}_i$ und $\mathbf{x}_j \cdot \mathbf{x}_j$,
die \textbf{Nebendiagonale} $-w_{ij}$ den \textbf{gemischten} Term
$-2\,\mathbf{x}_i \cdot \mathbf{x}_j$.
\end{block}
\end{columns}
\source{\cite{deLeeuw1977}}
\end{frame}


\begin{frame}{SMACOF: Zerlegung des Stress}
\footnotesize
\begin{columns}
\column{0.55\textwidth}
Mit der binomischen Formel lösen wir die Klammer in 
$\sigma(X) = \sum_{i<j} w_{ij}(\delta_{ij} - d_{ij}(X))^2$ auf:
\[
  \begin{aligned}
    \sigma(X)
      &= \sum_{i<j} w_{ij}
         \bigl(\delta_{ij}^2 - 2\,\delta_{ij}\,d_{ij}(X) + d_{ij}^2(X)\bigr)\\
      &= \sum_{i<j} w_{ij}\,\delta_{ij}^2
         - 2\sum_{i<j} w_{ij}\,\delta_{ij}\,d_{ij}(X)\\
      &\quad + \sum_{i<j} w_{ij}\,d_{ij}^2(X)
  \end{aligned}
\]
Die drei Summen sind der \textbf{konstante}, der \textbf{Problem}- und der
\textbf{quadratische} Term:
\[
  \sigma(X) = \underbrace{\sum_{i<j} w_{ij}\,\delta_{ij}^2}_{\text{konstant}}
  - 2\underbrace{\sum_{i<j} w_{ij}\,\delta_{ij}\,d_{ij}(X)}_{\text{Problemterm $\rho(X)$}}
  + \underbrace{\operatorname{tr}(X^T V X)}_{\text{quadratisch}}
\]

\column{0.40\textwidth}
\begin{blockC1}{\footnotesize Die drei Terme}
\begin{itemize}
  \item $\sum_{i<j} w_{ij}\,\delta_{ij}^2$: nur \textbf{Abstände im Ausgangsraum},
        \textbf{konstant}.
  \item $\sum_{i<j} w_{ij}\,\delta_{ij}\,d_{ij}(X)$: enthält $d_{ij}$,
        \textbf{nicht konvex}.
  \item $\sum_{i<j} w_{ij}\,d_{ij}^2(X) = \operatorname{tr}(X^T V X)$:
        \textbf{quadratisch}.
\end{itemize}
\end{blockC1}

\begin{blockC2}{\footnotesize Kern des Problems $\rho(X)$}
Nur der \textbf{Problemterm} ist schwierig. Genau hier setzt die
\textbf{Majorisierung} an.
\end{blockC2}
\end{columns}
\source{\cite{deLeeuw1977}}
\end{frame}

%%%%%%%%%%%%%%%%%%%%%%%%%%%%%%%%%%%%%%%%%%%%%
\begin{frame}{SMACOF: Die Idee der Majorante}
\footnotesize
\begin{columns}
\column{0.55\textwidth}
Wir minimieren den Stress \textbf{iterativ}. Sei $Z = X^{(k)}$ die
\textbf{aktuelle}, bereits bekannte Konfiguration des letzten Schritts.

\vspace{0.4em}
\begin{itemize}
  \item Unser schwieriger, nicht konvexer Term ist ja 
        $\rho(X) = \sum_{i<j} w_{ij}\,\delta_{ij}\,d_{ij}(X)$.
  \item Idee: Wir ersetzen $\rho(X)$ durch eine \textbf{lineare} untere
        Schranke, die von der bekannten Konfiguration $Z$ ausgeht.
  \item Da die Punkte $\mathbf{z}_i \in Z$ \textbf{fest} sind, lassen sich
        alle Abstände $d_{ij}(Z)$ konkret \textbf{berechnen}.
\end{itemize}

\column{0.40\textwidth}
\begin{blockC3}{\footnotesize Zwei Konfigurationen}
\begin{itemize}
  \item $Z$: \textbf{aktuelle} Konfiguration, bekannt und fest.
  \item $X$: \textbf{gesuchte} neue Konfiguration.
\end{itemize}
\end{blockC3}
\end{columns}
\source{\cite{deLeeuw1977}}
\end{frame}

\begin{frame}{SMACOF: Die Cauchy-Schwarz-Ungleichung}
\footnotesize
\begin{columns}
\column{0.55\textwidth}
Die \textbf{Cauchy-Schwarz}-Ungleichung besagt für zwei Vektoren
$\mathbf{a}, \mathbf{b}$:
\[
  \mathbf{a} \cdot \mathbf{b} \leq \lVert \mathbf{a} \rVert\,\lVert \mathbf{b} \rVert
\]
mit \textbf{Gleichheit} genau dann, wenn $\mathbf{a}$ und $\mathbf{b}$
\textbf{parallel} sind.

\vspace{0.4em}
Wir setzen $\mathbf{a} = \mathbf{x}_i - \mathbf{x}_j$ (neu) und
$\mathbf{b} = \mathbf{z}_i - \mathbf{z}_j$ (bekannt). Umgestellt ergibt das
eine \textbf{lineare} untere Schranke für den Abstand:
\[
  d_{ij}(X) \geq
  \frac{(\mathbf{x}_i - \mathbf{x}_j)\cdot(\mathbf{z}_i - \mathbf{z}_j)}
       {d_{ij}(Z)}
\]

\column{0.40\textwidth}
\begin{blockC3}{\footnotesize Der Trick}
Rechts steht $\mathbf{x}_i - \mathbf{x}_j$ nur noch \textbf{linear}: Der
Vektor $\mathbf{z}_i - \mathbf{z}_j$ und $d_{ij}(Z)$ sind \textbf{feste
Zahlen} aus $Z$. Die störende \textbf{Wurzel} ist verschwunden.
\end{blockC3}

\begin{blockC4}{\footnotesize Berührung}
Für $X = Z$ gilt \textbf{Gleichheit}.  Schranke und Stress stimmen am
aktuellen Punkt überein!
\end{blockC4}
\end{columns}
\source{\cite{deLeeuw1977}}
\end{frame}


\begin{frame}{SMACOF: Minimierung der Majorante}
\footnotesize
\begin{columns}
\column{0.55\textwidth}
Wendet man die Schranke auf alle Paare an, wird der Problemterm nach unten
beschränkt:
\[
  \rho(X) \geq \operatorname{tr}\bigl(X^T B(Z)\,Z\bigr)
\]
Setzt man das in den Stress ein, entsteht eine \textbf{quadratische
Majorante}. Ihr Minimum findet man durch Nullsetzen des Gradienten:
\[
  V X = B(Z)\,Z \;\Rightarrow\; X = V^{+} B(Z)\,Z
\]
$V$ ist \textbf{nicht invertierbar}: Alle Zeilensummen sind null, da eine
\textbf{Translation} aller Punkte die Abstände nicht ändert.

\column{0.40\textwidth}
\begin{block}{\footnotesize Matrix $B(Z)$}
\[
  b_{ij} =
  \begin{cases}
    -\dfrac{w_{ij}\,\delta_{ij}}{d_{ij}(Z)} & i \neq j,\; d_{ij}(Z) > 0\\[6pt]
    0 & i \neq j,\; d_{ij}(Z) = 0\\[4pt]
    -\sum_{k \neq i} b_{ik} & i = j
  \end{cases}
\]
\end{block}

\begin{blockDef}{\footnotesize Pseudoinverse $V^{+}$}
Die \textbf{Moore-Penrose}-Pseudoinverse verallgemeinert die Inverse für
\textbf{singuläre} Matrizen und liefert die im Ursprung \textbf{zentrierte}
Lösung.
\end{blockDef}
\end{columns}
\source{\cite{deLeeuw1977}}

\note{
\textbf{Singuläre Matrix:}
\begin{itemize}
    \item Eine singuläre Matrix ist eine quadratische Matrix, die \textbf{keine Inverse} besitzt.
    \item Äquivalent gilt: $\det(V)=0$.
\end{itemize}

\textbf{Moore-Penrose-Pseudoinverse:}
\begin{itemize}
    \item Die Pseudoinverse $V^{+}$ ist eine Verallgemeinerung der gewöhnlichen Inversen.
    \item Ist $V$ invertierbar, so gilt
    \[
        V^{+}=V^{-1}.
    \]
    \item Ist $V$ singulär, existiert keine gewöhnliche Inverse. Die Pseudoinverse liefert dennoch eine wohldefinierte Lösung.
\end{itemize}

\textbf{Anwendung auf lineare Gleichungssysteme:}
Für
\[
    Vx=b
\]
liefert
\[
    x=V^{+}b
\]
die Lösung mit \textbf{minimaler euklidischer Norm} unter allen möglichen Lösungen (bzw. die beste Least-Squares-Lösung, falls keine exakte Lösung existiert).

\textbf{Bezug zu SMACOF:}
Da die Matrix $V$ singulär ist (Zeilensummen gleich Null), kann $V^{-1}$ nicht berechnet werden. Daher wird die Pseudoinverse $V^{+}$ verwendet, die die eindeutig bestimmte, im Ursprung zentrierte Konfiguration liefert.
}
\end{frame}

\begin{frame}{SMACOF: Die Guttman-Transformation}
\footnotesize
\begin{columns}
\column{0.55\textwidth}
Für \textbf{gleiche} Gewichte ($w_{ij} = 1$) vereinfacht sich die
Pseudoinverse zu $V^{+} = \tfrac{1}{n}I$. Damit folgt die
\textbf{Guttman-Transformation}:
\[
  X^{(k+1)} = \frac{1}{n}\,B\!\left(X^{(k)}\right) X^{(k)}
\]

\begin{itemize}
  \item In jedem Schritt wird die aktuelle Konfiguration $X^{(k)}$ in eine
        \textbf{bessere} $X^{(k+1)}$ überführt.
  \item Der \textbf{Stress} sinkt dabei garantiert.
  \item Das Ergebnis $X^{(k+1)}$ wird zum neuen $Z$.
\end{itemize}

\column{0.40\textwidth}
\begin{blockC1}{\footnotesize Iteration}
Die Transformation wird \textbf{wiederholt} angewandt, bis sich der Stress
kaum noch ändert.
\end{blockC1}

\begin{blockC2}{\footnotesize Kern}
Ein schwieriges Minimierungsproblem wird so zu einer \textbf{Folge einfacher}
Matrixmultiplikationen.
\end{blockC2}
\end{columns}
\source{\cite{deLeeuw1977}}
\note{Eine Majorante ist eine Funktion, die überall oberhalb einer anderen Funktion liegt. Im SMACOF-Kontext ist sie das zentrale Werkzeug, um das schwierige Stress-Minimierungsproblem in eine Folge einfacher Schritte zu zerlegen.}
\end{frame}


%%%%%%%%%%%%%%%%%%%%%%%%%%%%%%%%%%%%%%%%%%%%%

\begin{frame}{SMACOF: Der Algorithmus}
\begin{columns}
\column{0.60\textwidth}
\footnotesize
\SetAlgoCaptionLayout{footnotesize}
\begin{algorithm}[H]
\SetAlgoLined
\KwIn{Unähnlichkeiten $\delta_{ij}$, Dimension $p$, Toleranz $\varepsilon$}
\KwOut{Konfiguration $X$}
Startkonfiguration $X^{(0)}$ initialisieren\;
$k = 0$\;
$\sigma_0 = \sigma(X^{(0)})$\;
\Repeat{$\sigma_{k-1} - \sigma_k < \varepsilon$ \textbf{oder} $k > k_{\max}$}{
  $k = k + 1$\;
  $B\!\left(X^{(k-1)}\right)$ berechnen\;
  \tcc{Guttman-Transformation}
  $X^{(k)} = \tfrac{1}{n} B\!\left(X^{(k-1)}\right) X^{(k-1)}$\; 
  $\sigma_k = \sigma(X^{(k)})$\;
}
\caption{SMACOF}
\end{algorithm}

\column{0.35\textwidth}
\begin{blockC1}{Ablauf}
In jedem Schritt wird die \textbf{Guttman-Transformation} angewandt und der
\textbf{Stress} neu berechnet bis er kaum noch sinkt.
\end{blockC1}
\end{columns}
\source{\cite{deLeeuw1977}}
\end{frame}

\begin{frame}{SMACOF: Eigenschaften}
\begin{columns}
\column{0.55\textwidth}
\begin{itemize}
  \item Der Stress ist über die Iterationen \textbf{monoton fallend}.
  \item Das Verfahren konvergiert gegen ein \textbf{lokales} Minimum.
  \item Das Ergebnis hängt von der \textbf{Startkonfiguration} ab.
  \item In der Praxis startet man oft mehrfach oder mit einer klassischen
        MDS-Lösung.
\end{itemize}

\column{0.40\textwidth}
\begin{blockC1}{Fazit}
SMACOF wandelt ein schwieriges Minimierungsproblem in eine \textbf{Folge
einfacher} quadratischer Schritte um. \textbf{Stabil} und garantiert
\textbf{konvergent}.
\end{blockC1}
\end{columns}
\source{\cite{deLeeuw1977}}
\end{frame}

\begin{frame}[fragile]{MDS in Python: Daten skalieren}
\begin{columns}
\column{0.55\textwidth}
MDS arbeitet mit \textbf{Abständen}. Die Daten müssen daher vorab
\textbf{skaliert} werden:
\begin{minted}[fontsize=\footnotesize]{python}
from sklearn.preprocessing import StandardScaler
scaler = StandardScaler()
X_transformed = scaler.fit_transform(X)
\end{minted}

\column{0.40\textwidth}
\begin{blockC1}{Warum skalieren?}
Ohne Skalierung würden Merkmale mit \textbf{großen Werten} die Abstände
\textbf{dominieren} und das Ergebnis verzerren.
\end{blockC1}
\end{columns}
\source{One Hot Fahrzeugdaten.ipynb}
\end{frame}

\begin{frame}[fragile]{MDS in Python: Projektion}
\begin{columns}
\column{0.47\textwidth}
Transformation aus $\mathbb{R}^7$ nach $\mathbb{R}^2$. \\ \medskip

Der Parameter
\texttt{n\_components} ist die \textbf{Dimension} des Zielraums. \\ \medskip

\texttt{fit\_transform()} macht training und transformation in einem Schritt. \\ \medskip


\begin{minted}[fontsize=\footnotesize]{python}
from sklearn.manifold import MDS
mds = MDS(
    n_components=2, 
    random_state=0)

X_2d = mds.fit_transform(X_transformed)
\end{minted}

\column{0.48\textwidth}
\centering
\includesvg[width=\linewidth]{figures/scatter_mds_auto_mpg.svg}

\end{columns}
\source{One Hot Fahrzeugdaten.ipynb}
\end{frame}


\begin{frame}[fragile]{MDS in Python: Projektion}
\begin{columns}
\column{0.47\textwidth}
\begin{itemize}
\item Man verwendet MDS oft, um sich einen \textbf{Überblick über Daten} zu verschaffen. 
\item Man kann so beispielsweise sehen, wie gut sich \textbf{Klassen separieren lassen}.
\item \textbf{Beispiel: }Projektion des 4-dimensionalen Iris-Datensatzes mit zugehörigen Klassen 
\end{itemize}

\column{0.48\textwidth}
\centering
\includesvg[width=\linewidth]{figures/scatter_mds_iris.svg}

\end{columns}
\source{MDS Iris mit Klassen.ipynb}
\end{frame}

\begin{frame}{Dimensionsreduktion: Grenzen von MDS}
\begin{columns}
\column{0.47\textwidth}
\begin{itemize}
  \item MDS eignet sich sehr gut für \textbf{kleine bis mittelgroße}
        Datensätze.
  \item Bei \textbf{großen} Datenmengen ist das Verfahren sehr
        \textbf{langsam}.
  \item Bei Daten \textbf{höherer Dimension} ist es sehr langsam und liefert oft \textbf{keine
        optimalen} Ergebnisse.
  \item \textbf{Beispiel:} die MDS-Projektion des \textbf{MNIST}-Datensatzes.
\end{itemize}

\column{0.48\textwidth}
\centering
\includesvg[width=\linewidth]{figures/scatter_mds_mnist.svg}
\end{columns}
\source{MDS MNIST.ipynb}
\end{frame}

\makequestionspage{\FRAGENIMAGE}



\begin{frame}{t-SNE: Grundidee}
\begin{columns}
\column{0.60\textwidth}
\begin{itemize}
  \item \textbf{t-SNE} (t-distributed Stochastic Neighbor Embedding) ist
        eine \textbf{nichtlineare} Dimensionsreduktion.
  \item Sie eignet sich besonders zur \textbf{Visualisierung}
        hochdimensionaler Daten in 2D oder 3D.
  \item \textbf{Idee:} Paarweise \textbf{Ähnlichkeiten} der Punkte werden als
        \textbf{Wahrscheinlichkeiten} ausgedrückt.
  \item Im Zielraum sucht man eine Anordnung, deren Ähnlichkeiten die des
        Ausgangsraums möglichst gut \textbf{nachbilden}.
\end{itemize}

\column{0.35\textwidth}
\begin{blockC1}{Fokus}
t-SNE erhält vor allem die \textbf{lokale} Nachbarschaft: nahe Punkte
bleiben nah, die globale Struktur ist zweitrangig.
\end{blockC1}
\end{columns}
\source{\cite{vanDerMaaten2008}}
\end{frame}


\begin{frame}{t-SNE: Ähnlichkeiten im Ausgangsraum}
\footnotesize
\begin{columns}
\column{0.47\textwidth}
\begin{itemize}
  \item Ausgehend von einem Punkt bestimmt man den \textbf{Abstand} zu allen
        anderen Punkten.
\end{itemize}

\vspace{0.3em}
\centering
\resizebox{0.9\linewidth}{!}{%
\begin{tikzpicture}[
  gr/.style={circle, fill=AccentColor, draw=AccentTitle,
             very thick, minimum size=6pt, inner sep=0pt},
  ma/.style={circle, fill=AlertColor, draw=AlertTitle,
             very thick, minimum size=6pt, inner sep=0pt},
  bl/.style={circle, fill=SecondaryColor, draw=SecondaryTitle,
             very thick, minimum size=6pt, inner sep=0pt}
]
% Achsen
\draw[thick, SecondaryTitle] (0,4.2) -- (0,0) -- (5.2,0);
% gruenes Cluster (oben links)
\node[gr] (g1) at (1.3,3.5) {};
\node[gr] (g2) at (1.0,2.9) {};
\node[gr] at (1.7,2.9) {};
\node[gr] at (1.3,2.4) {};
% Abstand-Markierung zwischen g1 und g2
\draw[ExampleTitle, thick] (g1) -- (g2);
\node[ExampleTitle, font=\scriptsize, rotate=61] at (0.8,3.3) {Abstand};
% blaues Cluster (rechts)
\node[bl] at (3.6,3.0) {};
\node[bl] at (3.3,2.5) {};
\node[bl] at (3.9,2.5) {};
\node[bl] at (3.6,2.0) {};
% magenta Cluster (unten)
\node[ma] at (1.7,1.3) {};
\node[ma] at (1.4,0.8) {};
\node[ma] at (2.0,0.8) {};
\node[ma] at (1.7,0.3) {};
\end{tikzpicture}}

\column{0.48\textwidth}
\begin{itemize}
  \item Um die \textbf{Nachbarschaftsstruktur} stärker zu gewichten, nimmt
        man die \textbf{Normalverteilung} mit dem aktuellen Punkt als Zentrum
        und berechnet daraus die \textbf{unskalierte Ähnlichkeit}.
\end{itemize}

\vspace{0.3em}
\centering
\resizebox{0.9\linewidth}{!}{%
\begin{tikzpicture}[
gr/.style={circle, fill=AccentColor, draw=AccentTitle,
             very thick, minimum size=6pt, inner sep=0pt}
]
% Gauss-Glocke (glatt, vorberechnet)
\draw[black, line width=1pt]
  (-3.6,0.001) -- (-3.15,0.005) -- (-2.7,0.028) -- (-2.25,0.11) --
  (-1.8,0.338) -- (-1.35,0.812) -- (-0.9,1.516) -- (-0.63,1.957) --
  (-0.45,2.206) -- (-0.27,2.39) -- (0,2.5) -- (0.27,2.39) --
  (0.45,2.206) -- (0.63,1.957) -- (0.9,1.516) -- (1.35,0.812) --
  (1.8,0.338) -- (2.25,0.11) -- (2.7,0.028) -- (3.15,0.005) -- (3.6,0.001);
% Kreuz auf der Kurve bei x=0.9 (Hoehe 1.516)
\node[cross out,
  draw=ExampleTitle,
  line width=1.5pt,
  minimum size=6pt,
  inner sep=0pt] at (0.9,1.516) {};
% Projektionslinie nach unten
\draw[ExampleTitle, thick, dashed] (0.9,1.516) -- (0.9,-0.7);
% Punktpaar unten
\node[gr] (p1) at (0.5,-0.9) {};
\node[gr] (p2) at (0.9,-0.9) {};
\draw[ExampleTitle, thick] (p1) -- (p2);
\end{tikzpicture}}
\end{columns}
\source{\cite{youtube_NEaUSP4YerM}}
\end{frame}


\begin{frame}{t-SNE: Ähnlichkeiten im Ausgangsraum}
\footnotesize
\begin{columns}
\column{0.47\textwidth}
\begin{itemize}
  \item Berechnet man die Ähnlichkeit zu einem \textbf{weiter entfernten}
        Punkt, zeigt sich der Sinn der \textbf{Normalverteilung}.
  \item \textbf{Weit entfernte} Punkte bekommen eine sehr \textbf{geringe}
        Ähnlichkeit, \textbf{nahe} Punkte eine \textbf{hohe}.
\end{itemize}

\vspace{0.4em}
\centering
\resizebox{0.9\linewidth}{!}{%
\begin{tikzpicture}[
 gr/.style={circle, fill=AccentColor, draw=AccentTitle,
             very thick, minimum size=6pt, inner sep=0pt},
  ma/.style={circle, fill=AlertColor, draw=AlertTitle,
             very thick, minimum size=6pt, inner sep=0pt}
]
% Gauss-Glocke (glatt, vorberechnet)
\draw[black, line width=1pt]
  (-3.6,0.001) -- (-3.15,0.005) -- (-2.7,0.028) -- (-2.25,0.11) --
  (-1.8,0.338) -- (-1.35,0.812) -- (-0.9,1.516) -- (-0.63,1.957) --
  (-0.45,2.206) -- (-0.27,2.39) -- (0,2.5) -- (0.27,2.39) --
  (0.45,2.206) -- (0.63,1.957) -- (0.9,1.516) -- (1.35,0.812) --
  (1.8,0.338) -- (2.25,0.11) -- (2.7,0.028) -- (3.15,0.005) -- (3.6,0.001);
% Kreuz am flachen Rand bei x=2.25 (Hoehe 0.11)
\node[cross out,
  draw=ExampleTitle,
  line width=1.5pt,
  minimum size=6pt,
  inner sep=0pt] at (2.25,0.11) {};
% Punktpaar unten, weit auseinander
\node[gr] (p1) at (-0.4,-0.5) {};
\node[ma] (p2) at (2.25,-0.5) {};
\draw[ExampleTitle, thick] (p1) -- (p2);
\draw[ExampleTitle, thick, dashed] (2.25,-0.5) -- (2.25,0.11);
\end{tikzpicture}}

\column{0.48\textwidth}
\centering
\resizebox{0.9\linewidth}{!}{%
\begin{tikzpicture}[
  gr/.style={circle, fill=AccentColor, draw=AccentTitle,
             very thick, minimum size=6pt, inner sep=0pt},
  ma/.style={circle, fill=AlertColor, draw=AlertTitle,
             very thick, minimum size=6pt, inner sep=0pt},
  bl/.style={circle, fill=SecondaryColor, draw=SecondaryTitle,
             very thick, minimum size=6pt, inner sep=0pt}
]
% Achsen
\draw[thick, SecondaryTitle] (0,4.2) -- (0,0) -- (5.2,0);
% gruenes Cluster (oben)
\node[gr] (g1) at (1.6,3.5) {};
\node[gr] (g2) at (1.3,2.9) {};
\node[gr] at (2.0,2.9) {};
\node[gr] at (1.7,2.4) {};
% blaues Cluster (rechts)
\node[bl] at (4.0,3.2) {};
\node[bl] at (3.7,2.7) {};
\node[bl] at (4.3,2.7) {};
\node[bl] at (4.0,2.2) {};
% magenta Cluster (unten)
\node[ma] (m1) at (1.5,1.3) {};
\node[ma] at (1.8,0.9) {};
\node[ma] at (2.1,1.2) {};
\node[ma] at (1.8,0.4) {};
% Verbindung gruen <-> magenta (weit entfernt)
\draw[ExampleTitle, thick] (g2) -- (m1);
\end{tikzpicture}}
\end{columns}
\source{\cite{youtube_NEaUSP4YerM}}
\end{frame}


\begin{frame}{t-SNE: Ähnlichkeiten im Ausgangsraum}
\footnotesize
\begin{columns}
\column{0.47\textwidth}
\begin{itemize}
  \item So berechnet man die \textbf{unskalierten} Unähnlichkeiten zu
        \textbf{allen} Datenpunkten.
  \item Im nächsten Teilschritt werden die Unähnlichkeiten so
        \textbf{skaliert}, dass sie sich für jeden Punkt auf \textbf{1}
        aufsummieren.
\end{itemize}

\vspace{0.4em}
\centering
\resizebox{0.9\linewidth}{!}{%
\begin{tikzpicture}[
gr/.style={circle, fill=AccentColor, draw=AccentTitle,
             very thick, minimum size=6pt, inner sep=0pt},
  ma/.style={circle, fill=AlertColor, draw=AlertTitle,
             very thick, minimum size=6pt, inner sep=0pt},
  bl/.style={circle, fill=SecondaryColor, draw=SecondaryTitle,
             very thick, minimum size=6pt, inner sep=0pt},
  kr/.style={
  cross out,
  draw=ExampleTitle,
  line width=1.5pt,
  minimum size=5pt,
  inner sep=0pt
}
]
% Gauss-Glocke (glatt, vorberechnet)
\draw[black, line width=1pt]
  (-3.6,0.001) -- (-3.15,0.005) -- (-2.7,0.028) -- (-2.25,0.11) --
  (-1.8,0.338) -- (-1.35,0.812) -- (-0.9,1.516) -- (-0.63,1.957) --
  (-0.45,2.206) -- (-0.27,2.39) -- (0,2.5) -- (0.27,2.39) --
  (0.45,2.206) -- (0.63,1.957) -- (0.9,1.516) -- (1.35,0.812) --
  (1.8,0.338) -- (2.25,0.11) -- (2.7,0.028) -- (3.15,0.005) -- (3.6,0.001);
% Kreuze magenta (weit links, flach)
\foreach \x in {-3.0,-2.7,-2.5} \node[kr] at (\x,0.05)  {};
\node[kr] at (-2.2,0.15)  {};
% Kreuze gruen (nah, oben rechts der Spitze)
\node[kr] at (0.2,2.4)  {};
\node[kr] at (0.5,2.1)  {};
\node[kr] at (0.6,1.95)  {};
% Kreuze blau (weit rechts, flach)
\foreach \x in {2.7,2.9, 3.1} \node[kr] at (\x,0.05)  {};
\node[kr] at (2.4,0.1)  {};
% Punktgruppen unter der Achse
\foreach \x in {-3.0,-2.7, -2.5,-2.2} \node[ma] at (\x,-0.55) {};
\foreach \x in {0.0,0.2,0.5, 0.6} \node[gr] at (\x,-0.55) {};
\foreach \x in {2.4,2.7,2.9,3.1} \node[bl] at (\x,-0.55) {};
\end{tikzpicture}}

\column{0.48\textwidth}
\centering
\resizebox{0.9\linewidth}{!}{%
\begin{tikzpicture}[
  gr/.style={circle, fill=AccentColor, draw=AccentTitle,
             very thick, minimum size=6pt, inner sep=0pt},
  ma/.style={circle, fill=AlertColor, draw=AlertTitle,
             very thick, minimum size=6pt, inner sep=0pt},
  bl/.style={circle, fill=SecondaryColor, draw=SecondaryTitle,
             very thick, minimum size=6pt, inner sep=0pt}
]
% Achsen
\draw[thick, SecondaryTitle] (0,4.2) -- (0,0) -- (5.2,0);
% gruenes Cluster (oben, Bezugspunkt g0)
\node[gr] (g0) at (1.4,3.0) {};
\node[gr] (g1) at (1.7,3.6) {};
\node[gr] (g2) at (2.0,3.0) {};
\node[gr] (g3) at (1.7,2.5) {};
% blaues Cluster (rechts)
\node[bl] (b1) at (4.0,3.3) {};
\node[bl] (b2) at (3.7,2.8) {};
\node[bl] (b3) at (4.3,2.8) {};
\node[bl] (b4) at (4.0,2.3) {};
% magenta Cluster (unten)
\node[ma] (m1) at (1.5,1.3) {};
\node[ma] (m2) at (1.9,0.9) {};
\node[ma] (m3) at (2.2,1.3) {};
\node[ma] (m4) at (1.8,0.4) {};
% Verbindungen vom Bezugspunkt g0 zu allen anderen
\foreach \t in {b1,b2,b3,b4,m1,m2,m3,m4,g1,g2,g3}
  \draw[ExampleTitle, thick, opacity=0.7] (g0) -- (\t);
\end{tikzpicture}}
\end{columns}
\source{\cite{youtube_NEaUSP4YerM}}
\end{frame}

\begin{frame}{t-SNE: Ähnlichkeiten im Ausgangsraum}
\begin{columns}
\column{0.47\textwidth}
\textbf{Warum skaliert man die Ähnlichkeiten?}

\vspace{0.4em}
\begin{itemize}
  \item Zusammengehörige Bereiche müssen nicht zwingend gleich \textbf{dicht}
        besiedelt sein.
  \item Auch wenn der \textbf{nächste Nachbar} etwas weiter entfernt ist,
        bleibt er der \textbf{nächste Nachbar}.
  \item Das wird durch die \textbf{Skalierung} berücksichtigt.
\end{itemize}

\column{0.48\textwidth}
\centering
\resizebox{0.9\linewidth}{!}{%
\begin{tikzpicture}[
gr/.style={circle, fill=AccentColor, draw=AccentTitle,
             very thick, minimum size=6pt, inner sep=0pt},
  ma/.style={circle, fill=AlertColor, draw=AlertTitle,
             very thick, minimum size=6pt, inner sep=0pt},
  bl/.style={circle, fill=SecondaryColor, draw=SecondaryTitle,
             very thick, minimum size=6pt, inner sep=0pt}
]
% Achsen
\draw[thick, SecondaryTitle] (0,4.5) -- (0,0) -- (6.0,0);
% gruenes Cluster (oben links, dicht)
\node[gr] at (1.3,4.0) {};
\node[gr] at (1.0,3.5) {};
\node[gr] at (1.6,3.5) {};
\node[gr] at (1.3,3.0) {};
% magenta Cluster (unten links, dicht)
\node[ma] at (1.5,2.1) {};
\node[ma] at (1.2,1.6) {};
\node[ma] at (1.8,1.6) {};
\node[ma] at (1.5,1.1) {};
% blaues Cluster (rechts, weit auseinandergezogen)
\node[bl] at (4.6,3.0) {};
\node[bl] at (4.0,1.9) {};
\node[bl] at (5.5,1.9) {};
\node[bl] at (4.6,0.9) {};
\end{tikzpicture}}
\end{columns}
\source{\cite{youtube_NEaUSP4YerM}}
\end{frame}


\begin{frame}{t-SNE: Ähnlichkeiten im Ausgangsraum}
\footnotesize
\begin{columns}
\column{0.55\textwidth}
Mit unseren Vorbetrachtungen lässt sich die Ähnlichkeit von $\mathbf{x}_j$ zu $\mathbf{x}_i$ als \textbf{bedingte
Wahrscheinlichkeit} definieren:
\[
  p_{j|i} =
  \frac{\exp\!\bigl(-\lVert \mathbf{x}_i - \mathbf{x}_j \rVert^2 / 2\sigma_i^2\bigr)}
       {\sum_{k \neq i} \exp\!\bigl(-\lVert \mathbf{x}_i - \mathbf{x}_k \rVert^2 / 2\sigma_i^2\bigr)}
\]

\begin{itemize}
  \item $p_{j|i}$ ist \textbf{groß} für nahe, \textbf{klein} für ferne
        Punkte.
  \item Die Breite $\sigma_i$ wird pro Punkt über die \textbf{Perplexität}
        festgelegt.
  \item In dichten Regionen wird $\sigma_i$ \textbf{klein}, in dünnen
        \textbf{groß}.
\end{itemize}

\column{0.40\textwidth}
\begin{blockC1}{\footnotesize Perplexität}
Ein Maß für die \textbf{effektive Anzahl} der Nachbarn (typisch $5$--$50$).
Sie steuert das $\sigma_i$ jedes Punktes.
\end{blockC1}
\end{columns}
\source{\cite{vanDerMaaten2008}}
\end{frame}

\begin{frame}{t-SNE: Symmetrisierung}
\footnotesize
\begin{columns}
\column{0.55\textwidth}
Die bedingten Wahrscheinlichkeiten sind \textbf{nicht symmetrisch}
($p_{j|i} \neq p_{i|j}$). Man symmetrisiert sie zu einer gemeinsamen
Verteilung:
\[
  p_{ij} = \frac{p_{j|i} + p_{i|j}}{2n}
\]

\begin{itemize}
  \item $n$ ist die Anzahl der Datenpunkte.
  \item Es gilt $p_{ij} = p_{ji}$ und $\sum_{i,j} p_{ij} = 1$.
  \item Die $p_{ij}$ bilden die \textbf{Ähnlichkeitsmatrix} $P$ des
        Ausgangsraums.
\end{itemize}

\column{0.40\textwidth}
\centering
\textbf{Matrix $P$}\\[3pt]
\resizebox{0.9\linewidth}{!}{%
\begin{tikzpicture}
\foreach \i in {0,...,4} {
  \foreach \j in {0,...,4} {
    \pgfmathsetmacro{\val}{\i==\j ? 0 : 0.3 + 0.5*rnd}
    \pgfmathsetmacro{\dval}{\i==\j ? 1 : \val}
    \fill[SecondaryColor, opacity=\dval] (\j,-\i) rectangle ++(1,-1);
    \draw[white, thin] (\j,-\i) rectangle ++(1,-1);
  }
}
\end{tikzpicture}}

{\scriptsize Dunkel $=$ hohe Ähnlichkeit}
\end{columns}
\source{\cite{vanDerMaaten2008}}
\end{frame}

\begin{frame}{t-SNE: Ähnlichkeiten im Zielraum}
\footnotesize
\begin{columns}
\column{0.55\textwidth}

Die Ähnlichkeitsverteiung im Zielraum wird im Prinzip gleich berechnet. Allerdings nutzt t-SNE hier eine
\textbf{Student-t}-Verteilung mit einem Freiheitsgrad (Cauchy):
\[
  q_{ij} =
  \frac{\bigl(1 + \lVert \mathbf{y}_i - \mathbf{y}_j \rVert^2\bigr)^{-1}}
       {\sum_{k \neq l} \bigl(1 + \lVert \mathbf{y}_k - \mathbf{y}_l \rVert^2\bigr)^{-1}}
\]

\begin{itemize}
  \item Die $q_{ij}$ bilden die \textbf{Ähnlichkeitsmatrix} $Q$ des
        Zielraums.
  \item Ziel: $Q$ soll $P$ möglichst gut \textbf{nachbilden}.
  \item Der Nenner \textbf{normiert} über alle Paare.
\end{itemize}

\column{0.40\textwidth}
\centering
\textbf{Matrix $Q$}\\[3pt]
\resizebox{0.9\linewidth}{!}{%
\begin{tikzpicture}
\foreach \i in {0,...,4} {
  \foreach \j in {0,...,4} {
    \pgfmathsetmacro{\val}{\i==\j ? 0 : 0.25 + 0.55*rnd}
    \pgfmathsetmacro{\dval}{\i==\j ? 1 : \val}
    \fill[ExampleColor, opacity=\dval] (\j,-\i) rectangle ++(1,-1);
    \draw[white, thin] (\j,-\i) rectangle ++(1,-1);
  }
}
\end{tikzpicture}}

{\scriptsize Ziel: $Q \approx P$}
\end{columns}
\source{\cite{vanDerMaaten2008}}
\end{frame}

\begin{frame}{t-SNE: Warum die Student-t-Verteilung?}
\footnotesize
\begin{columns}
\column{0.47\textwidth}
\begin{itemize}
  \item In hohen Dimensionen ist \textbf{viel mehr Platz} als in 2D: ferne
        Punkte drängen sich beim Projizieren zusammen (\textbf{Crowding
        Problem}).
  \item Die Student-t-Verteilung hat \textbf{schwerere Ränder} als die
        Gauß-Glocke.
  \item Dadurch dürfen \textbf{unähnliche} Punkte im Zielraum \textbf{weiter}
        auseinander liegen.
  \item Cluster werden so klar \textbf{getrennt} dargestellt.
\end{itemize}

\column{0.48\textwidth}
\centering
\includesvg[width=\linewidth]{figures/vergleich_gauss_student.svg}

{\scriptsize Schwere Ränder $\Rightarrow$ mehr Platz für ferne Punkte}
\end{columns}
\source{\cite{vanDerMaaten2008}}
\end{frame}

\begin{frame}{t-SNE: Kostenfunktion}
\footnotesize
\begin{columns}
\column{0.55\textwidth}
Die Anordnung im Zielraum wird so gewählt, dass $Q$ möglichst nah an $P$
liegt. Als Maß dient die \textbf{Kullback-Leibler-Divergenz}:
\[
  C = \mathrm{KL}(P \,\Vert\, Q)
    = \sum_{i \neq j} p_{ij} \log \frac{p_{ij}}{q_{ij}}
\]

\begin{itemize}
  \item Sie misst, wie stark sich $Q$ von $P$ \textbf{unterscheidet}.
  \item Sie ist \textbf{asymmetrisch}: nahe Paare ($p_{ij}$ groß) im Zielraum
        auseinanderzureißen wird \textbf{stark bestraft}.
  \item Ferne Paare zusammenzuziehen kostet dagegen \textbf{wenig}.
\end{itemize}

\column{0.40\textwidth}
\begin{blockC1}{\footnotesize Folge}
t-SNE erhält bevorzugt die \textbf{lokale} Struktur: enge Nachbarschaften
werden treu wiedergegeben.
\end{blockC1}
\end{columns}
\source{\cite{vanDerMaaten2008}}
\end{frame}

\begin{frame}{t-SNE: Optimierung}
\footnotesize
\begin{columns}
\column{0.55\textwidth}
Die Kostenfunktion wird per \textbf{Gradientenabstieg} minimiert. Der
Gradient hat eine anschauliche Form:
\[
  \frac{\partial C}{\partial \mathbf{y}_i}
  = 4 \sum_{j} (p_{ij} - q_{ij})
    \bigl(1 + \lVert \mathbf{y}_i - \mathbf{y}_j \rVert^2\bigr)^{-1}
    (\mathbf{y}_i - \mathbf{y}_j)
\]

\begin{itemize}
  \item Der Term $(p_{ij} - q_{ij})$ wirkt wie eine \textbf{Kraft} zwischen
        den Punkten.
  \item $p_{ij} > q_{ij}$: \textbf{Anziehung}; $p_{ij} < q_{ij}$:
        \textbf{Abstoßung}.
  \item Das Verfahren pendelt sich in ein \textbf{Gleichgewicht} ein.
\end{itemize}

\column{0.40\textwidth}
\begin{blockC2}{Hinweise}
\begin{itemize}
  \item Die Kostenfunktion ist \textbf{nicht konvex}. Das Ergebnis hängt
        vom \textbf{Zufallsstart} ab.
  \item \textbf{Perplexität} und Iterationszahl beeinflussen das Bild stark.
\end{itemize}
\end{blockC2}
\end{columns}
\source{\cite{vanDerMaaten2008}}
\end{frame}

\begin{frame}{t-SNE: Eigenschaften und Praxis}
\begin{columns}
\column{0.55\textwidth}
\begin{itemize}
  \item t-SNE trennt \textbf{Cluster} oft sehr überzeugend.
  \item \textbf{Abstände} zwischen Clustern und deren \textbf{Größen} sind
        \textbf{nicht} zuverlässig interpretierbar.
  \item Das Verfahren ist \textbf{rechenintensiv} und für sehr große
        Datenmengen langsam.
  \item Es dient der \textbf{Visualisierung}, nicht als Vorverarbeitung für
        andere Modelle.
\end{itemize}

\column{0.40\textwidth}
\begin{blockC1}{Merke}
t-SNE zeigt \textbf{welche} Punkte zusammengehören, nicht \textbf{wie weit}
Gruppen voneinander entfernt sind.
\end{blockC1}
\end{columns}
\source{\cite{vanDerMaaten2008}}
\end{frame}


\begin{frame}[fragile]{t-SNE in Python}
\begin{columns}
\column{0.60\textwidth}
\centering
\includesvg[width=\linewidth]{figures/tsne_mnist.svg}

\column{0.40\textwidth}
\begin{minted}[fontsize=\footnotesize]{python}
from sklearn.manifold import TSNE

tsne = TSNE(n_components=2,
            random_state=0,
            perplexity=50)
X_2d = tsne.fit_transform(X)
\end{minted}

\vspace{0.4em}
\footnotesize
\begin{itemize}
  \item \texttt{n\_components}: \textbf{Zieldimension} (hier $2$).
  \item \texttt{perplexity}: effektive \textbf{Nachbarzahl}.
\end{itemize}
\end{columns}
\source{t-SNE MNIST.ipynb}
\end{frame}

\makequizpage{\QUIZURL}{\QUIZQRCODE}{\QUIZIMAGE}

\begin{frame}{Übung: t-SNE auf dem Iris-Datensatz}
\begin{columns}
\column{0.60\textwidth}
Führen Sie auf dem \textbf{Iris-Datensatz} eine Dimensionsreduktion mit
\textbf{t-SNE} durch.

\vspace{0.4em}
\begin{itemize}
  \item Variieren Sie dabei die \textbf{Perplexität}.
  \item Geben Sie den \textbf{projizierten} Datensatz aus.
\end{itemize}

\column{0.35\textwidth}
\centering
\twemoji[width=0.9\linewidth]{laptop}
\end{columns}
\source{Lösung: t-SNE.ipynb}
\end{frame}


\makequestionspage{\FRAGENIMAGE}